%% This is file `StellarexPoster.tex',
%% it is part of the Stellarex LaTeX bundle
%% (derived from the IPP bundle by Marei Peischl, LPPL 1.3c)
%%
%% Generic poster template showing the box styles, figures and the bibliography.
%% Compile with: latexmk StellarexPoster   (lualatex + biber, see .latexmkrc)
\documentclass[
	fontsize=26pt,% base font size (24-28pt is readable from 1.5-2 m); title and box headings scale with it
	paperwidth=48in,% passed to geometry
	paperheight=36in,% passed to geometry
%	margin=20mm,% width of the navy frame
%	logoheight=60mm,% height of the Stellarex logo in the title block
%	footerheight=30mm,% minimum height of the footer area
]{stellarex-poster}

%language setup
\usepackage[english]{babel}
\usepackage{microtype}
\usepackage[autostyle]{csquotes}
\usepackage{amsmath}
\usepackage{booktabs}
\usepackage{lipsum}% placeholder text, remove when done

% Figures and tables: use \captionof inside poster boxes (floats are not possible)
\graphicspath{{}{figures/}{logos/}}

\usepackage[style=numeric-comp]{biblatex}
\addbibresource{StellarexPoster.bib}

\begin{document}

\title{Poster Title Goes Here, Ideally on a Single Line}

% Authors: \inst{} refers to the institute list, \contact{} adds a corresponding author to the footer
\author{First Author\inst{1}\contact{first.author@stellarex.com} \and Second Author\inst{1,2} \and Third Author\inst{2}}
\institute{\inst{1}Stellarex, Inc., Toronto, Canada \and\inst{2}Partner Institution, City, Country}

\footer{Name of the Conference, City, Date}
% Affiliation logo(s) next to the Stellarex logo, top right (placeholder: PPPL)
% height defaults to logoheight; change it with \StellarexPosterSetup{affiliationlogoheight=40mm}
\setAffiliationLogos{pppl-logo}
% Additional logo(s) below the Stellarex logo in the title block
% \setHeadlineLogos{partner-logo.pdf}
% Logos in the footer, e.g. of partners or funding agencies
\footergraphics{\includegraphics[height=20mm]{example-image}\hspace{1em}}
% Funding statement, bottom right
\acknowledgement{This work was supported by Funding Agency, grant number 0000000.}

\begin{tcbposter}[
	poster={
		columns=4,% number of equal width columns
		rows=4,% number of rows
	}]

%%% Column 1 %%%%%%%%%%%%%%%%%%%%%%%%%%%%%%%%%%%%%%%%%%%%%%%%%%%%%%%%%%%%%%%%%%%
\begin{posterboxenv}[
	title=Abstract,
	abstract,% navy title bar
	]{name=abstract, column=1, row=1}
	The \texttt{abstract} style has a navy title bar. Use it for the abstract or summary at the top left of the poster.

	\lipsum[1][1-4]
\end{posterboxenv}

\begin{posterboxenv}[title=Introduction]{name=intro, column=1, between=abstract and bottom}
	Boxes without a style option use the default style: a grey title bar and a white body.
	Text can be structured with lists and equations as usual:
	\begin{itemize}
		\item First point, with a citation~\cite{placeholder-article}
		\item Second point, with a reference to Eq.~\eqref{eq:example}
		\item Third point, citing a book~\cite{placeholder-book}
	\end{itemize}
	\begin{equation}
		\label{eq:example}
		\nabla \cdot \mathbf{B} = 0, \qquad \mathbf{J} \times \mathbf{B} = \nabla p
	\end{equation}
	\lipsum[2][1-3]

	\section{Section heading inside a box}
	Use \verb+\section+ to add headings inside a box. They are styled like box titles.

	\lipsum[3][1-4]

	\section{Another heading}
	\lipsum[4][1-4]
\end{posterboxenv}

%%% Column 2 %%%%%%%%%%%%%%%%%%%%%%%%%%%%%%%%%%%%%%%%%%%%%%%%%%%%%%%%%%%%%%%%%%%
\begin{posterboxenv}[title=Methods]{name=methods, column=2, row=1, rowspan=2}
	A single figure at full box width, with a numbered caption:

	\medskip
	\centerline{\includegraphics[width=\linewidth,height=.5\linewidth]{example-image}}
	\captionof{figure}{A placeholder figure. Replace \texttt{example-image} with your own file, e.g.\ \texttt{figures/my-plot.pdf}.\label{fig:methods}}

	Figures are referenced as usual (Fig.~\ref{fig:methods}). The method follows~\cite{placeholder-proceedings}:
	\begin{enumerate}
		\item First step of the method
		\item Second step of the method
		\item Third step of the method
	\end{enumerate}
	\lipsum[11][1-5]
\end{posterboxenv}

%%% Column 3 %%%%%%%%%%%%%%%%%%%%%%%%%%%%%%%%%%%%%%%%%%%%%%%%%%%%%%%%%%%%%%%%%%%
\begin{posterboxenv}[title=Results]{name=results, column=3, row=1, rowspan=2}
	Two figures side by side, each with its own caption:

	\medskip
	\begin{minipage}[t]{.48\linewidth}
		\includegraphics[width=\linewidth,height=.75\linewidth]{example-image-a}
		\captionof{figure}{Left placeholder figure.\label{fig:left}}
	\end{minipage}\hfill
	\begin{minipage}[t]{.48\linewidth}
		\includegraphics[width=\linewidth,height=.75\linewidth]{example-image-b}
		\captionof{figure}{Right placeholder figure.\label{fig:right}}
	\end{minipage}

	\bigskip
	Tables work the same way (Table~\ref{tab:results}):

	\medskip
	\begin{center}
		\captionof{table}{A placeholder table.\label{tab:results}}
		\begin{tabular}{lrr}
			\toprule
			Case & Quantity A & Quantity B \\
			\midrule
			Baseline & 1.00 & 0.00 \\
			Variant 1 & 1.25 & 0.10 \\
			Variant 2 & 1.50 & 0.20 \\
			\bottomrule
		\end{tabular}
	\end{center}

	\lipsum[6][1-4]
\end{posterboxenv}

%%% Columns 2 and 3, bottom rows %%%%%%%%%%%%%%%%%%%%%%%%%%%%%%%%%%%%%%%%%%%%%%%
\begin{posterboxenv}[title=Discussion]{name=discussion, column=2, span=2, below=methods}
	This box spans two columns (\verb+span=2+), which suits wide figures.

	\medskip
	\begin{minipage}[t]{.6\linewidth}
		\vspace{0pt}% aligns the tops of the two minipages
		\includegraphics[width=\linewidth,height=.3\linewidth]{example-image-16x9}
		\captionof{figure}{A wide placeholder figure.\label{fig:wide}}
	\end{minipage}\hfill
	\begin{minipage}[t]{.36\linewidth}
		\vspace{0pt}
		\lipsum[14][1-5]
	\end{minipage}
\end{posterboxenv}

\begin{posterboxenv}[
	title=Key result,
	highlight,% gold callout (default color of the highlight style)
	]{name=key, column=2, span=2, between=discussion and bottom}
	\begin{minipage}[c]{.6\linewidth}
		The \texttt{highlight} style draws attention to the main message. Without an argument it uses Stellarex gold.

		\lipsum[8][1-3]
	\end{minipage}\hfill
	\begin{minipage}[c]{.36\linewidth}
		\includegraphics[width=\linewidth,height=.5\linewidth]{example-image-c}
		\captionof{figure}{A figure next to text.}
	\end{minipage}
\end{posterboxenv}

%%% Column 4 %%%%%%%%%%%%%%%%%%%%%%%%%%%%%%%%%%%%%%%%%%%%%%%%%%%%%%%%%%%%%%%%%%%
\begin{posterboxenv}[
	title=Conclusions,
	conclusion,% navy title bar, light navy body
	]{name=conclusion, column=4, row=1}
	The \texttt{conclusion} style has a navy title bar and a tinted body.
	\lipsum[15][1-2]
	\begin{itemize}
		\item \lipsum[9][1]
		\item \lipsum[9][3]
		\item \lipsum[9][5]
	\end{itemize}
\end{posterboxenv}

\begin{posterboxenv}[
	title=Outlook,
	highlight=blue,% other colors: navy, blue_light, orange, grey, grey_light, ...
	]{name=outlook, column=4, below=conclusion}
	Use \verb+highlight=<color>+ for any palette color.
	\begin{itemize}
		\item \lipsum[10][1]
		\item \lipsum[10][3]
	\end{itemize}
\end{posterboxenv}

\begin{posterboxenv}[title=References]{name=references, column=4, between=outlook and bottom}
	\nocite{*}% lists all entries of the .bib file, remove when done
	\printbibliography[heading=none]
\end{posterboxenv}

\end{tcbposter}

\end{document}
